%!TEX root = main.tex


\subsection{Correctness of ExpandWhens}
We have formalized the following central statement of the ExpandWhens transformation:

\begin{quote}
For every ground-type writable (sub)component,
ExpandWhens retains the last connect or \prettyll{is invalid} statement
that appears under any combination of \prettyll{when} conditions in the module.
\linebreak
If the (sub)component is declared under multiple conditions,
ExpandWhens selects between them using a multiplexer or \prettyll{validif} expression.
\end{quote}

To describe the dependency on (a combination of) conditions,
we introduce an auxiliary notion:
\begin{definition}
An \emph{expression tree} is a binary tree whose internal nodes are annotated with conditions
(usually, conditions found in \prettyll{when} statements)
and whose leaves are annotated with expressions or one of three special constants:
{\tt invalid}, {\tt not connected}, or {\tt undeclared}.
\end{definition}
An expression tree is able to describe the declarations and connects of one (sub)component under all conditions.

This tree serves as the basis to define a function {\tt expand\_branch\_one\_var\_sem}
that constructs an expression tree for one ground-type (sub)component.
It takes as parameters a statement sequence $ss$, a ground-type (sub)component $\mathit{ref}$, and a default expression tree.
It constructs the expression tree for $\mathit{ref}$, based on the conditions and connect statements in $ss$.
If $ss$ does not contain a connect under every condition, the default value is inserted in these case(s).

In the next step we construct {\tt expand\_branch\_sem\_conform},
a predicate that, given the statement sequence of a module or of a \prettyll{when} branch,
checks for all writable components
whether the implementation expression of {\tt expand\_branch} corresponds to the semantic tree of {\tt expand\_branch\_\linebreak[0]one\_var\_\linebreak[0]sem}.
The function receives type and writability information from two component environments:
one for all components in the current module,
and one for the ports of all other modules (that may be accessed in instances of these modules).

The proof requires mutual induction over statement sequences and statements:
if statement sequence $ss$ and statement $s$ satisfy the correspondence,
then also statement sequence $ss ; s$.
A statement \prettyll{when} $c$ \prettyll{:} $ss^\mathsf{T}$ \prettyll{else} $ss^\mathsf{F}$
satisfies the correspondence if both $ss^\mathsf{T}$ and $ss^\mathsf{F}$ satisfy it.

%To find (sub)components of ground type,
%one needs type and writability information about the components,
%which we receive from the component environment.

To guarantee that the type and writability information can actually be retrieved from the component environments,
we need a one-to-one correspondence between a component environment and a connect map:
there is an entry in the connect map
iff there is a writable component declaration in the component environment.
We therefore proved a lemma stating that {\tt expand\_branch}
preserves the correspondence.
Again we had to use mutual induction over statement sequences and statements.

This mutual induction also required us to define the two Coq datatypes mutually inductive;
we cannot use the sequence or list types from the standard library.